\documentclass[12pt,twoside]{article}
%\usepackage[table]{xcolor} % important to avoid options clash.
%\input{02465shared_preamble}
%\usepackage{cleveref}
\usepackage{url}
\usepackage{graphics}
\usepackage{multicol}
\usepackage{rotate}
\usepackage{rotating}
\usepackage{booktabs}
\usepackage{hyperref}
\usepackage{pifont}
\usepackage{latexsym}
\usepackage[english]{babel}
\usepackage{epstopdf}
\usepackage{etoolbox}
\usepackage{amsmath}
\usepackage{amssymb}
\usepackage{multirow,epstopdf}
\usepackage{fancyhdr}
\usepackage{booktabs}
\usepackage{xcolor}
\newcommand\redt[1]{ {\textcolor[rgb]{0.60, 0.00, 0.00}{\textbf{ #1} } } }


\newcommand{\m}[1]{\boldsymbol{ #1}}
\newcommand{\yoursolution}{ \redt{(your solution here) } } 



\title{ Report 1 hand-in }
\date{ \today }
\author{Alice (\texttt{s000001})\and  Bob (\texttt{s000002})\and Clara (\texttt{s000003}) } 

\begin{document}
\maketitle

\begin{table}[ht!]
\caption{Attribution table. Feel free to add/remove rows and columns}
\begin{tabular}{llll}
\toprule
                                                        & Alice   & Bob    & Clara   \\
\midrule
 1: A basic blaster-business                            & 0-100\%  & 0-100\% & 0-100\%  \\
 2: Warmup                                              & 0-100\%  & 0-100\% & 0-100\%  \\
 3: Manually computing $J_{N-1}$                        & 0-100\%  & 0-100\% & 0-100\%  \\
 4: Compute optimal policy and value function           & 0-100\%  & 0-100\% & 0-100\%  \\
 5: Kiosk2                                              & 0-100\%  & 0-100\% & 0-100\%  \\
 6: Policy explanation continued                        & 0-100\%  & 0-100\% & 0-100\%  \\
 7: Go east                                             & 0-100\%  & 0-100\% & 0-100\%  \\
 8: Predict consequence of actions                      & 0-100\%  & 0-100\% & 0-100\%  \\
 9: Possible future states                              & 0-100\%  & 0-100\% & 0-100\%  \\
 10: Shortest path                                      & 0-100\%  & 0-100\% & 0-100\%  \\
 11: Predict consequence of actions with one ghost      & 0-100\%  & 0-100\% & 0-100\%  \\
 12: Possible future states with one ghost              & 0-100\%  & 0-100\% & 0-100\%  \\
 13: Optimal one-ghost planning                         & 0-100\%  & 0-100\% & 0-100\%  \\
 14: Predict consequence of actions with several ghosts & 0-100\%  & 0-100\% & 0-100\%  \\
 15: Future states                                      & 0-100\%  & 0-100\% & 0-100\%  \\
 16: Optimal planning                                   & 0-100\%  & 0-100\% & 0-100\%  \\
\bottomrule
\end{tabular}
\end{table}

%\paragraph{Statement about collaboration:}
%Please edit this section to reflect how you have used external resources. The following statement will in most cases suffice: 
%\emph{The code in the irls/project1 directory is entirely}

%\paragraph{Main report:}
Headings have been inserted in the document for readability. You only have to edit the part which says \yoursolution. 

\section{The kiosk (\texttt{kiosk.py})}
\subsubsection*{{\color{red}Problem 1:  A basic blaster-business}}

\yoursolution 	
\redt{Include the following elements: \begin{align}
	N & = 14 \\
	\mbox{for $k=0,\dots,N$: }\quad	\mathcal{S}_k & = \{0,1, \dots, 20\\} \\
	\mbox{for $k=0,\dots,N-1$: }\quad \mathcal{A}_k(x_k) & = \{0,1, \dots, 15\\} \\
    \mbox{for $k=0,\dots,N-1$: }\quad f_k(x_k, u_k, w_k) & = max(0, min(n_s, x + u - w)) \\
	\mbox{for $k=0,\dots,N-1$: }\quad g_k(x_k, u_k, w_k) & = c_o \cdot u_k - s_p \cdot min(x_k + u_k, w_k) \\
	g_N(x_N) & = 0 \\
    \mbox{for $k=0,\dots,N-1$: }\quad P_W(w_k | x_k, u_k) & = {0:.3, 3:.6, 6:0.1}
\end{align} }

\subsubsection*{{\color{red}Problem 3:  Manually computing $J_{N-1}$}}
		
	\yoursolution The last step of the DP algorithm gives:
	$$
	J_{N-1}(20)  = \min_{u_{N-1} } \EE_{w_{N-1} }\left[ g_{N-1}(20, u_{N-1}, w_{N-1} ) +  g_N(x_N) \right]
	$$
	Therefore:
	As the the final cost (g_n) is zero that term can be ignored. Moreover, the cost function g_{N-1} can be inserted:
	J_{N-1}(20) = \min_{u_{N-1} } \EE_{w_{N-1} }\left[ 1.5 \cdot u_{N-1} - 2.1 \cdot min(20 + u_{N-1}, w_{N-1})\right]
	As we are told the stock is 20, the minimum of 20 + u_{N-1} and w_{N-1} will always be w_{N-1}. 
	Therefore, the cost function can be simplified to:
	J_{N-1}(20) = \min_{u_{N-1} } \EE_{w_{N-1} }\left[ 1.5 \cdot u_{N-1} - 2.1 \cdot w_{N-1}\right]
	Because the expectation is linear, the expectation can be put inside the brackets as follows:
	J_{N-1}(20) = \min_{u_{N-1} } \left[ 1.5 \cdot u_{N-1} - 2.1 \cdot \EE_{w_{N-1} }\right]
	Then the expectation of w can be calculated, but as it does not depend on u, we can minimise 1.5*u with the smallest possible action u=0, resulting in 1.5*0=0.
	Therefore, the cost function can be simplified to:
	J_{N-1}(20) = 0 - 2.1 \cdot (0 \cdot 0.3 + 3 \cdot 0.6 + 6 \cdot 0.1) = -5.04
	Thereby the expected profit when starting with fulls tock and choosing action u=0 is 5.04.
	
\subsubsection*{{\color{red}Problem 6:  Policy explanation continued}}
	
	$$\mu_{N-1}(0) = \argmin_u \EE_{w} \left[g_k(0, u, w) + g_N( f(0,u,w) ) \right] $$
\yoursolution

Computing how many blasters would be bought the last day given that the inventory is empty, 
when following the first policy
$$\mu_{N-1}(0) = \argmin_u \EE_{w} \left[g_k(0, u, w) + g_N( f(0,u,w) ) \right] $$
$$\mu_{N-1}(0) = \argmin_u \EE_{w} \left[c_o*u - s_p*min(0+u,w) + 0 \right]  $$
$$\mu_{N-1}(0) = \argmin_u \EE_{w} \left[c_o*u - s_p*min(0+u,w) + 0 \right] $$

Now calcualting the expected value for the demand:
$$\mu_{N-1}(0) = \argmin_u \left[c_o*u - s_p*\EE_{w}[min(0+u,w)] + 0 \right] $$
$$ \EE_{w}[min(0+u,w)] = 0.3*min(0+u,0) + 0.6*min(0+u,3) + 0.1*min(0+u,6) $$

Since W belongs to {0,3,6}, the E[min(u,W)] is a piecewise linear function of u, with breakpoints at 0,3 and 6.
This means that the minimum of the cost function will be at one of those breakpoints.
Therefore, we can evaluate the cost function at those breakpoints to find the minimum cost:
u=0: J(0) = 1.5*0 - 2.1*(0.3*0 + 0.6*0 + 0.1*0) = 0
u=3: J(3) = 1.5*3 - 2.1*(0.3*0 + 0.6*3 + 0.1*3) = 0.09
u=6: J(6) = 1.5*6 - 2.1*(0.3*0 + 0.6*3 + 0.1*6) = 3.96

This shows that the smallet cost is obtained, when u=0.
This means that according to the first policy, the optimal action is to not buy any blasters the last day.

\section{Avoid the droid (\texttt{pacman.py)}} 
\end{document}